% GENERATED FILE. Edit canonical lessons, then run npm run generate:books.
% canonical-source-hash: fnv1a64:57e4f7dae25ab263
% canonical-lessons: BN-C468-hritpindo, BN-C468-jokrit, BN-C468-pakostholi, BN-C468-shira, BN-C468-dhukpuk-kora

\chapter{Inside the Body}
\label{ch:bn-inside-the-body}
\hlchaptermodality{468}

\begin{chapteropening}
\textbf{By the end of this chapter:} \emph{I can name the heart, the liver, the stomach and the veins.}

Complete the last lesson of chapter 468: I can name the heart, the liver, the stomach and the veins.
\end{chapteropening}

\section[\texorpdfstring{\bn{হৃৎপিণ্ড} (hritpinḍô)}{হৃৎপিণ্ড (hritpinḍô)}]{\bn{হৃৎপিণ্ড} (hritpinḍô) — the heart (the organ)}

\label{lesson:BN-C468-hritpindo}



\begin{quote}
Before the new one: say the Bengali for the voice, then the Bengali for to knit one's brows.
\end{quote}

\subsection*{You'll want to know: \bn{হৃৎপিণ্ড}}
\textbf{\bn{হৃৎপিণ্ড}} — \emph{hritpinḍô} — \textquotedblleft{}the heart (the organ)\textquotedblright{}.

\begin{quote}
\bn{ব্যায়াম} \bn{করলে} \bn{হৃৎপিণ্ড} \bn{ভালো} \bn{থাকে।}
\end{quote}

\begin{quote}
\emph{bæyām kôrle hritpinḍô bhālo thāke.}
\end{quote}

\begin{quote}
Exercise keeps the heart healthy.
\end{quote}

\bn{হৃদয়} (taught) is the heart of feeling; \bn{হৃৎপিণ্ড} is the pump the doctor listens to.

\begin{cousinweb}
\textbf{\bn{হৃৎপিণ্ড}} (\emph{hritpinḍô}) is a tatsama, Sanskrit \emph{hṛtpiṇḍa}, \textquotedblleft{}heart-mass\textquotedblright{}, the heart (Monier-Williams).
\end{cousinweb}

\begin{grammarlens}[title={What you've built}]
One more word for describing what you lived and felt.
\end{grammarlens}

\subsection*{Your turn}
\begin{itemize}
\raggedright
  \item \emph{Say it:} \emph{hritpinḍô}
  \item \emph{Say it:} \emph{hritpinḍô}, once more
  \item \emph{Say it:} \emph{bæyām kôrle hritpinḍô bhālo thāke.} — Exercise keeps the heart healthy.
  \item \emph{Recall:} say the Bengali for the voice, then the Bengali for to knit one's brows, then say \emph{hritpinḍô} again
\end{itemize}

\begin{tcolorbox}[breakable,colback=teal!4,colframe=teal!35!black,title={Before you move on}]
What does \bn{হৃৎপিণ্ড} mean? (\textquotedblleft{}The heart (the organ)\textquotedblright{}.) Say it once more.
\end{tcolorbox}
\section[\texorpdfstring{\bn{যকৃৎ} (jôkrit)}{যকৃৎ (jôkrit)}]{\bn{যকৃৎ} (jôkrit) — the liver}

\label{lesson:BN-C468-jokrit}



\begin{quote}
Before the new one: say the Bengali for to knit one's brows, then the Bengali for the heart (the organ).
\end{quote}

\subsection*{You'll want to know: \bn{যকৃৎ}}
\textbf{\bn{যকৃৎ}} — \emph{jôkrit} — \textquotedblleft{}the liver\textquotedblright{}.

\begin{quote}
\bn{দাদুর} \bn{যকৃৎ} \bn{দুর্বল।}
\end{quote}

\begin{quote}
\emph{dādur jôkrit durbôl.}
\end{quote}

\begin{quote}
Grandfather's liver is weak.
\end{quote}

\begin{cousinweb}
\textbf{\bn{যকৃৎ}} (\emph{jôkrit}) is a tatsama, Sanskrit \emph{yakṛt}, \textquotedblleft{}the liver\textquotedblright{} (Monier-Williams), which the dictionary compares with Latin \emph{jecur}.
\end{cousinweb}

\begin{grammarlens}[title={What you've built}]
One more word for describing what you lived and felt.
\end{grammarlens}

\subsection*{Your turn}
\begin{itemize}
\raggedright
  \item \emph{Say it:} \emph{jôkrit}
  \item \emph{Say it:} \emph{jôkrit}, once more
  \item \emph{Say it:} \emph{dādur jôkrit durbôl.} — Grandfather's liver is weak.
  \item \emph{Recall:} say the Bengali for to knit one's brows, then the Bengali for the heart (the organ), then say \emph{jôkrit} again
\end{itemize}

\begin{tcolorbox}[breakable,colback=teal!4,colframe=teal!35!black,title={Before you move on}]
What does \bn{যকৃৎ} mean? (\textquotedblleft{}The liver\textquotedblright{}.) Say it once more.
\end{tcolorbox}
\section[\texorpdfstring{\bn{পাকস্থলী} (pākôsthôlī)}{পাকস্থলী (pākôsthôlī)}]{\bn{পাকস্থলী} (pākôsthôlī) — the stomach (the organ)}

\label{lesson:BN-C468-pakostholi}



\begin{quote}
Before the new one: say the Bengali for the heart (the organ), then the Bengali for the liver.
\end{quote}

\subsection*{You'll want to know: \bn{পাকস্থলী}}
\textbf{\bn{পাকস্থলী}} — \emph{pākôsthôlī} — \textquotedblleft{}the stomach (the organ)\textquotedblright{}.

\begin{quote}
\bn{খাবার} \bn{প্রথমে} \bn{পাকস্থলীতে} \bn{যায়।}
\end{quote}

\begin{quote}
\emph{khābār prôthôme pākôsthôlīte jāy.}
\end{quote}

\begin{quote}
Food goes first to the stomach.
\end{quote}

\begin{cousinweb}
\textbf{\bn{পাকস্থলী}} (\emph{pākôsthôlī}) puts Sanskrit \emph{pāka}, \textquotedblleft{}cooking\textquotedblright{}, here digesting, with \emph{sthalī}, \textquotedblleft{}place, ground\textquotedblright{} (Monier-Williams): the place where food is cooked down.
\end{cousinweb}

\begin{grammarlens}[title={What you've built}]
One more word for describing what you lived and felt.
\end{grammarlens}

\subsection*{Your turn}
\begin{itemize}
\raggedright
  \item \emph{Say it:} \emph{pākôsthôlī}
  \item \emph{Say it:} \emph{pākôsthôlī}, once more
  \item \emph{Say it:} \emph{khābār prôthôme pākôsthôlīte jāy.} — Food goes first to the stomach.
  \item \emph{Recall:} say the Bengali for the heart (the organ), then the Bengali for the liver, then say \emph{pākôsthôlī} again
\end{itemize}

\begin{tcolorbox}[breakable,colback=teal!4,colframe=teal!35!black,title={Before you move on}]
What does \bn{পাকস্থলী} mean? (\textquotedblleft{}The stomach (the organ)\textquotedblright{}.) Say it once more.
\end{tcolorbox}
\section[\texorpdfstring{\bn{শিরা} (shirā)}{শিরা (shirā)}]{\bn{শিরা} (shirā) — a vein}

\label{lesson:BN-C468-shira}



\begin{quote}
Before the new one: say the Bengali for the liver, then the Bengali for the stomach (the organ).
\end{quote}

\subsection*{You'll want to know: \bn{শিরা}}
\textbf{\bn{শিরা}} — \emph{shirā} — \textquotedblleft{}a vein\textquotedblright{}.

\begin{quote}
\bn{নার্সটি} \bn{হাতের} \bn{শিরায়} \bn{ইনজেকশন} \bn{দিলেন।}
\end{quote}

\begin{quote}
\emph{nārsṭi hāter shirāy injekshôn dilen.}
\end{quote}

\begin{quote}
The nurse gave the injection into a vein in the arm.
\end{quote}

\begin{cousinweb}
\textbf{\bn{শিরা}} (\emph{shirā}) is a tatsama, Sanskrit \emph{sirā} (also written \emph{śirā}), \textquotedblleft{}a vessel of the body, a vein\textquotedblright{} (Monier-Williams).
\end{cousinweb}

\begin{grammarlens}[title={What you've built}]
One more word for describing what you lived and felt.
\end{grammarlens}

\subsection*{Your turn}
\begin{itemize}
\raggedright
  \item \emph{Say it:} \emph{shirā}
  \item \emph{Say it:} \emph{shirā}, once more
  \item \emph{Say it:} \emph{nārsṭi hāter shirāy injekshôn dilen.} — The nurse gave the injection into a vein in the arm.
  \item \emph{Recall:} say the Bengali for the liver, then the Bengali for the stomach (the organ), then say \emph{shirā} again
\end{itemize}

\begin{tcolorbox}[breakable,colback=teal!4,colframe=teal!35!black,title={Before you move on}]
What does \bn{শিরা} mean? (\textquotedblleft{}A vein\textquotedblright{}.) Say it once more.
\end{tcolorbox}
\section[\texorpdfstring{\bn{ধুকপুক} \bn{করা} (dhukpuk kôrā)}{ধুকপুক করা (dhukpuk kôrā)}]{\bn{ধুকপুক} \bn{করা} (dhukpuk kôrā) — to thump; to throb (of the heart)}

\label{lesson:BN-C468-dhukpuk-kora}



\begin{quote}
Before the new one: say the Bengali for the stomach (the organ), then the Bengali for a vein.
\end{quote}

\subsection*{You'll want to know: \bn{ধুকপুক} \bn{করা}}
\textbf{\bn{ধুকপুক} \bn{করা}} — \emph{dhukpuk kôrā} — \textquotedblleft{}to thump; to throb (of the heart)\textquotedblright{}.

\begin{quote}
\bn{ফলাফলের} \bn{আগে} \bn{মীরার} \bn{বুক} \bn{ধুকপুক} \bn{করছিল।}
\end{quote}

\begin{quote}
\emph{phôlāphôler āge mīrār buk dhukpuk kôrchhilô.}
\end{quote}

\begin{quote}
Before the results Mira's heart was thumping.
\end{quote}

Past (he/she …): \bn{সেই} \bn{ছাত্রীর} \bn{বুক} \bn{ধুকপুক} \bn{করল}, \emph{sei chhātrīr buk dhukpuk kôrlô}.

\begin{cousinweb}
\textbf{\bn{ধুকপুক} \bn{করা}} (\emph{dhukpuk kôrā}) is \emph{dhukpuk}, an imitative word for the beat of the heart, with \emph{kôrā}, \textquotedblleft{}to do\textquotedblright{} (CDIAL 2814).
\end{cousinweb}

\begin{grammarlens}[title={What you've built}]
One more word for describing what you lived and felt.
\end{grammarlens}

\subsection*{Your turn}
\begin{itemize}
\raggedright
  \item \emph{Say it:} \emph{dhukpuk kôrā}
  \item \emph{Say it:} \emph{dhukpuk kôrā}, once more
  \item \emph{Say it:} \emph{phôlāphôler āge mīrār buk dhukpuk kôrchhilô.} — Before the results Mira's heart was thumping.
  \item \emph{Recall:} say the Bengali for the stomach (the organ), then the Bengali for a vein, then say \emph{dhukpuk kôrā} again
\end{itemize}

\begin{tcolorbox}[breakable,colback=teal!4,colframe=teal!35!black,title={Before you move on}]
What does \bn{ধুকপুক} \bn{করা} mean? (\textquotedblleft{}To thump; to throb (of the heart)\textquotedblright{}.) Say it once more.
\end{tcolorbox}
